\begin{textsample}{POS Dim 1 – vem\_ed\_06 – Score 12.00 – vem\_ed\_06/t023.txt}  \label{ex:f1_pos_127}
Caminhos com Viseu Realizar \textbf{intercâmbios} virtuais nas Fatecs, empregando a língua portuguesa.
Com essa \textbf{perspectiva} em mente, Osvaldo Succi Jr., coordenador dos Projetos Colaborativos Internacionais (PCIs), fez uma série de conversas ao \textbf{longo} de 2020 com representantes do Instituto Politécnico de Viseu (IPV), em Portugal, para \textbf{iniciar} PCIs com as Fatecs.
A instituição pública portuguesa oferece cursos técnicos de graduação e pós-graduação, em diversas áreas, como saúde, educação, agrária, tecnologia e gestão.
Em 22 de fevereiro de 2021, \textbf{foi} realizada uma reunião na plataforma Teams, para delinear caminhos \textbf{possíveis} para as colaborações luso-brasileiras.
Do IPV, participaram Susana Amante (professora de Língua Inglesa), António Figueiredo (atuante na área de Gestão, no Departamento de Madeiras), Joaquim Antunes (com intervenção em áreas de Marketing e Turismo), Pedro Reis (Finanças), Jorge Manuel Martins (Diretor do Departamento de Madeiras) e João Luís Pereira.
Do lado brasileiro, esteve presente a equipe dos PCIs, além de Ricardo Sérgio Neiva Nóbrega (coordenador do curso de Comércio Exterior da Fatec Indaiatuba) e Carlos Augusto Amaral Moreira (professor da Fatec Americana).
Um primeiro passo \textbf{proposto} pelos professores \textbf{é} a criação de projetos-piloto para, por exemplo, aferir \textbf{diferenças} culturais em \textbf{ambientes} de negócio no Brasil e em Portugal.
Após esse \textbf{encontro} virtual, Susana Amante, interlocutora da aproximação entre as instituições, concedeu o seguinte depoimento: Os projetos \textbf{colaborativos} internacionais constituem uma mais-valia para uma aprendizagem \textbf{ativa}, num \textbf{processo} construtivo em \textbf{que} os vários intervenientes – neste caso do IPV, em Portugal, e de Fatecs do Centro Paula Souza, no Brasil – se articulam para \textbf{que}, no âmbito das suas Unidades Curriculares, se envolvam em aprendizagens em \textbf{que} refletem acerca das suas próprias \textbf{perspectivas} e experiências, bem como sobre a \textbf{forma} como colegas, com outras culturas e \textbf{formas} de pensar, percepcionam o mundo \textbf{que} os rodeia e como compreendem e aplicam os \textbf{conteúdos} trabalhados.
Muito mais do \textbf{que} uma ênfase no individual, o \textbf{foco} \textbf{é} no coletivo.
A aprendizagem \textbf{colaborativa} implica um olhar sobre os \textbf{processos} para a resolução de desafios em \textbf{que} \textbf{cada} \textbf{elemento} contribui para uma equipa \textbf{interdisciplinar} e, muitas vezes, transdisciplinar, com o objetivo \textbf{comum} de resolução de problemas \textbf{cada} vez mais globais, porque \textbf{somos} \textbf{todos} cidadãos do mundo.
Neste sentido, o IPV congratula-se com esta parceria e espera \textbf{que} estes primeiros projetos-piloto venham a dar lugar a muitos outros PCIs produtivos, em anos vindouros.

% matched lemmas: ambiente, ativo, cada, colaborativo, comum, conteúdo, diferença, elemento, encontro, foco, forma, iniciar, intercâmbio, interdisciplinar, longo, perspectiva, possível, processo, propor, que, ser, todo
\end{textsample}
